% ============================================================
%  Section 6.3 – Hardware Integration
% ============================================================
\section{Hardware Integration}
\label{sec:ch6_hardware}

This section describes the physical assembly, wiring, and network configuration of the
collaborative robotic cell.  The goal is to document every connection that must be present for
the software stack described in the preceding sections to operate correctly.

% ------------------------------------------------------------------
\subsection{Cell Layout and Component Inventory}
\label{subsec:ch6_cell_layout}
% ------------------------------------------------------------------

The robotic cell is built around a shared assembly table to which both manipulators are rigidly
mounted.  Table~\ref{tab:ch6_hw_inventory} lists every hardware component and its role.

\begin{table}[H]
\centering
\caption{Hardware component inventory of the collaborative robotic cell.}
\label{tab:ch6_hw_inventory}
\setlength{\tabcolsep}{5pt}
\begin{tabular}{lll}
\toprule
\textbf{Component}         & \textbf{Role}                          & \textbf{Interface} \\
\midrule
ABB IRB~120                & Primary industrial manipulator         & ABB IRC5 / RWS (Ethernet) \\
AR4 Mk3 Research Arm       & Secondary collaborative arm            & Teensy 4.1 (\texttt{/dev/ttyACM0}) \\
ABB IRC5 Controller        & Low-level ABB motion controller        & Ethernet (192.168.125.1) \\
Intel RealSense D435i       & Global scene perception                & USB~3.0 \\
Station Laptop             & Perception \& executor node            & ZeroTier / Wi-Fi \\
Control Laptop             & Planner, MoveIt, Digital Twin          & ZeroTier / Wi-Fi \\
Teensy 4.1                 & AR4 joint-angle and gripper controller & USB-CDC (\texttt{/dev/ttyACM0}) \\
ABB Gripper (custom)       & Primary end-effector                   & ABB I/O board \\
AR4 Gripper (custom)       & Secondary end-effector                 & Teensy GPIO \\
\bottomrule
\end{tabular}
\end{table}

Figure~\ref{fig:ch6_cell_diagram} provides a schematic top-view of the shared worktable,
indicating the placement of both arms, the camera overhead mount, and the brick assembly region.

% ------------------------------------------------------------------
%  Figure: cell layout top-view (TikZ)
% ------------------------------------------------------------------
\begin{figure}[H]
\centering
\begin{tikzpicture}[
    font=\small\sffamily,
    scale=1.0,
    table/.style={draw, thick, fill=gray!10, minimum width=9cm, minimum height=5cm},
    arm/.style={draw, thick, fill=blue!20, rounded corners=4pt,
                minimum width=1.4cm, minimum height=1.4cm, align=center},
    camera/.style={draw, thick, fill=yellow!30, circle, minimum size=0.8cm, align=center},
    brick/.style={draw, thick, fill=red!20, rounded corners=2pt,
                  minimum width=2.5cm, minimum height=1.2cm, align=center},
    lbl/.style={font=\footnotesize\sffamily}
]

% Table surface
\draw[draw=gray!60, fill=gray!10, rounded corners=5pt] (-4.8,-2.8) rectangle (4.8,2.8);
\node[lbl, gray] at (0,-2.4) {Shared Assembly Table};

% ABB arm (left side)
\node[arm] (abb) at (-3.2, 0) {\textbf{ABB}\\IRB~120};
\draw[thick, ->, gray] (abb.east) -- ++(0.6,0);

% AR4 arm (right side)
\node[arm] (ar4) at (3.2, 0)  {\textbf{AR4}\\Mk3};
\draw[thick, ->, gray] (ar4.west) -- ++(-0.6,0);

% Assembly/handover zone (centre)
\node[brick] (zone) at (0, 0.6) {Assembly Zone};
\node[brick, fill=orange!20] (hzone) at (0,-0.6) {Handover Zone};

% Camera overhead
\node[camera] (cam) at (0, 2.1) {\scriptsize RS};
\node[lbl, above=0.1cm of cam] {RealSense D435i (overhead)};
\draw[dashed, gray] (cam.south) -- (zone.north);

% Workspace boundary arcs
\draw[dashed, blue!60, thick] (abb.center) circle (2.2cm);
\draw[dashed, orange!70!black, thick] (ar4.center) circle (2.2cm);

% Legend
\node[lbl, blue!70] at (-3.2,-2.1)  {ABB workspace};
\node[lbl, orange!70!black] at (3.2,-2.1) {AR4 workspace};

\end{tikzpicture}
\caption{Schematic top-view of the collaborative robotic cell.
         Dashed circles indicate approximate reachable workspace for each arm.
         The hatched overlap at the centre contains the handover zone.}
\label{fig:ch6_cell_diagram}
\end{figure}

% ------------------------------------------------------------------
\subsection{Wiring and Physical Connections}
\label{subsec:ch6_wiring}
% ------------------------------------------------------------------

Figure~\ref{fig:ch6_wiring} shows every electrical connection in the cell.  All inter-system
communication is either Ethernet, USB, or Wi-Fi; no proprietary fieldbuses are used, which
simplifies integration and debugging.

% ------------------------------------------------------------------
%  Figure: wiring diagram (TikZ)
% ------------------------------------------------------------------
\begin{figure}[H]
    \centering
    \resizebox{\textwidth}{!}{%
    \begin{tikzpicture}[
        font=\scriptsize,
        node distance=0.7cm and 1.4cm,
        hw/.style={draw, thick, rounded corners=2pt, align=center, fill=blue!10,
                   minimum width=2.6cm, minimum height=0.85cm},
        sensor/.style={draw, thick, rounded corners=2pt, align=center, fill=red!12,
                       minimum width=2.6cm, minimum height=0.85cm},
        laptop/.style={draw, rounded corners=2pt, align=center,
                       minimum width=3.0cm, minimum height=1.0cm},
        station/.style={laptop, fill=orange!12},
        ctrl/.style={laptop, fill=green!12},
        cloud/.style={draw, rounded corners=10pt, align=center, fill=cyan!12,
                      minimum width=2.8cm, minimum height=1.0cm},
        eth/.style={-Latex, thick, color=blue!70!black},
        prop/.style={-Latex, thick, color=gray!70!black},
        usb/.style={-Latex, thick, color=purple!70!black},
        cam/.style={-Latex, thick, color=red!70!black},
        wifi/.style={Latex-Latex, thick, color=cyan!60!black},
        lbl/.style={font=\tiny, fill=white, inner sep=1.5pt}
    ]
        % Hardware column (left)
        \node[hw]                    (irc5) {ABB IRC5\\controller};
        \node[hw, below=of irc5]     (abb)  {ABB IRB 120};
        \node[hw, below=of abb]      (ar4)  {AR4 Mk3\\(Teensy 4.1)};

        % Station laptop (centre, aligned with abb)
        \node[station, right=2.0cm of abb] (stat) {Station laptop};

        % Camera above station
        \node[sensor, above=1.1cm of stat] (camera) {RealSense D435i};

        % ZeroTier overlay (right of station)
        \node[cloud, right=of stat]  (zt)   {ZeroTier\\overlay};

        % Control laptop (far right)
        \node[ctrl, right=of zt]     (ctrl) {Control laptop};

        % Connections
        \draw[eth]  (irc5.east) -- node[lbl, above, sloped]{Ethernet 192.168.125.0/24} (stat.north west);
        \draw[prop] (abb.east)  -- node[lbl, above]{proprietary cable} (stat.west);
        \draw[usb]  (ar4.east)  -- node[lbl, below, sloped]{USB-CDC \texttt{/dev/ttyACM0}} (stat.south west);
        \draw[cam]  (camera.south) -- node[lbl, right]{USB 3.0} (stat.north);
        \draw[wifi] (stat.east) -- node[lbl, above]{Wi-Fi} (zt.west);
        \draw[wifi] (zt.east)   -- node[lbl, above]{Wi-Fi} (ctrl.west);

        \begin{scope}[on background layer]
            \node[draw, dashed, rounded corners=4pt,
                  fit=(irc5)(abb)(ar4), inner sep=0.25cm,
                  label=above:Cell hardware] {};
            \node[draw, dashed, rounded corners=4pt,
                  fit=(camera)(stat), inner sep=0.25cm,
                  label=below:Local node] {};
            \node[draw, dashed, rounded corners=4pt,
                  fit=(zt)(ctrl), inner sep=0.25cm,
                  label=above:Remote network] {};
        \end{scope}
    \end{tikzpicture}%
    }
    \caption{Physical and logical wiring of the collaborative robotic cell. Blue links are Ethernet, purple links are USB, the red link is USB 3.0 for the camera, gray is the proprietary ABB cable, and cyan double-arrows are Wi-Fi through the ZeroTier overlay.}
    \label{fig:ch6_wiring}
\end{figure}

\textbf{ABB IRC5 connection.} The station laptop communicates with the IRC5 controller over a 1\,Gbps Ethernet link using the
\textbf{ABB Robot Web Services (RWS)} HTTP/WebSocket interface.  The controller is assigned the
static address \texttt{192.168.125.1}; the station laptop is assigned \texttt{192.168.125.100}
on the same dedicated switch.  The \texttt{abb\_rws\_client} ROS~2 package translates RWS
responses into standard \texttt{/joint\_states} messages.

\textbf{AR4 Mk3 connection.} The AR4 arm is controlled by a Teensy~4.1 microcontroller mounted on the arm base.  The
microcontroller enumerates as a USB CDC serial device on the station laptop at
\texttt{/dev/ttyACM0}.  The ROS~2 AR4 driver serialises joint targets as
\texttt{GOTO\_JOINTS} commands, receives encoder feedback, and publishes joint states at 50\,Hz.

\textbf{Intel RealSense D435i.} The depth camera is mounted directly above the assembly table on a rigid aluminium bracket.
It connects to the station laptop via USB~3.0.  The \texttt{realsense2\_camera} ROS~2 driver
launches with aligned depth and colour streams; a YOLO detection node subscribes to the colour
stream and publishes brick poses on \texttt{/yolo/detections}.

% ------------------------------------------------------------------
\subsection{Network Configuration}
\label{subsec:ch6_network_config}
% ------------------------------------------------------------------

Three independent IP address spaces are active in the cell simultaneously
(Table~\ref{tab:ch6_ip}).  Routing between them is not required: each space is confined to a
single point-to-point link.

\begin{table}[H]
\centering
\caption{IP address assignments for the three network interfaces in the cell.}
\label{tab:ch6_ip}
\setlength{\tabcolsep}{5pt}
\begin{tabular}{llll}
\toprule
\textbf{Interface}          & \textbf{Subnet}       & \textbf{Station}    & \textbf{Control / Device} \\
\midrule
Ethernet (ABB)              & 192.168.125.0/24      & 192.168.125.100     & IRC5: 192.168.125.1 \\
ZeroTier virtual NIC        & 10.19.44.0/24         & 10.19.44.30         & 10.19.44.52 \\
Wi-Fi (physical, outbound)  & ISP-assigned (private)& DHCP                & DHCP \\
\bottomrule
\end{tabular}
\end{table}

On startup, each laptop runs the following initialisation sequence before any ROS~2 nodes are
launched:

\begin{enumerate}
    \item Verify ZeroTier membership and that both virtual IPs are reachable
          (\texttt{ping 10.19.44.52} from the station; \texttt{ping 10.19.44.30} from the
          control laptop).

    \item Set \texttt{RMW\_IMPLEMENTATION=rmw\_fastrtps\_cpp} and
          \texttt{FASTRTPS\_DEFAULT\_PROFILES\_FILE} to the shared XML profile, then export
          \texttt{ROS\_DOMAIN\_ID}.

    \item On the station laptop, confirm the IRC5 is reachable
          (\texttt{ping 192.168.125.1}) and the AR4 Teensy has enumerated
          (\texttt{ls /dev/ttyACM*}).

    \item Start Syncthing as a background daemon and confirm the shared
          \texttt{program\_storage} folder shows \emph{Up to Date} on both devices.
\end{enumerate}

Figure~\ref{fig:ch6_startup_sequence} formalises this boot order as a flowchart.

% ------------------------------------------------------------------
%  Figure: startup / initialisation flowchart (TikZ)
% ------------------------------------------------------------------
\begin{figure}[H]
    \centering
    \resizebox{0.72\textwidth}{!}{%
    \begin{tikzpicture}[
        font=\scriptsize,
        node distance=0.55cm and 1.1cm,
        start/.style={draw, rounded corners=8pt, align=center, fill=green!18,
                      minimum width=4.2cm, minimum height=0.85cm},
        finish/.style={draw, rounded corners=8pt, align=center, fill=green!35,
                       minimum width=4.2cm, minimum height=0.85cm},
        proc/.style={draw, rounded corners=2pt, align=center, fill=blue!10,
                     minimum width=4.2cm, minimum height=0.85cm},
        dec/.style={draw, diamond, aspect=2.1, align=center, inner sep=2pt,
                    minimum width=3.6cm, minimum height=0.9cm, fill=yellow!18},
        retry/.style={draw, rounded corners=2pt, align=center, fill=red!10,
                      minimum width=3.2cm, minimum height=0.85cm},
        arrow/.style={-Latex, thick},
        lbl/.style={font=\tiny, fill=white, inner sep=1.5pt}
    ]
        \node[start]                  (s1) {Power on both laptops};
        \node[proc,  below=of s1]     (p1) {Start ZeroTier daemon};
        \node[dec,   below=of p1]     (d1) {ZeroTier peers reachable?};
        \node[proc,  below=of d1]     (p2) {Export ROS env vars};
        \node[dec,   below=of p2]     (d2) {IRC5 \& Teensy available?};
        \node[proc,  below=of d2]     (p3) {Start Syncthing daemon};
        \node[dec,   below=of p3]     (d3) {Sync folder Up to Date?};
        \node[proc,  below=of d3]     (p4) {Launch ROS 2 nodes\\\texttt{ros2 launch master\_launch.py}};
        \node[finish, below=of p4]    (fin) {Cell ready};

        \node[retry, right=of d1] (r1) {Check Wi-Fi};
        \node[retry, right=of d2] (r2) {Check USB / Ethernet};
        \node[retry, right=of d3] (r3) {Wait for sync};

        % Main vertical Yes path
        \foreach \a/\b in {s1/p1, p1/d1, d1/p2, p2/d2, d2/p3, p3/d3, d3/p4, p4/fin}
            \draw[arrow] (\a) -- (\b);
        \node[lbl, right=2pt] at ($(d1.south)!0.5!(p2.north)$) {Yes};
        \node[lbl, right=2pt] at ($(d2.south)!0.5!(p3.north)$) {Yes};
        \node[lbl, right=2pt] at ($(d3.south)!0.5!(p4.north)$) {Yes};

        % No / retry: parallel offset arrows for each decision
        \foreach \d/\r in {d1/r1, d2/r2, d3/r3} {
            \draw[arrow] ([yshift=-0.18cm]\d.east) -- node[lbl, below]{No}
                         ([yshift=-0.18cm]\r.west);
            \draw[arrow] ([yshift=0.18cm]\r.west) -- node[lbl, above]{retry}
                         ([yshift=0.18cm]\d.east);
        }

        \begin{scope}[on background layer]
            \node[draw, dashed, rounded corners=4pt,
                  fit=(d1)(r1)(p2), inner sep=0.18cm,
                  label=right:Network check] {};
            \node[draw, dashed, rounded corners=4pt,
                  fit=(d2)(r2)(p3), inner sep=0.18cm,
                  label=right:Hardware check] {};
            \node[draw, dashed, rounded corners=4pt,
                  fit=(d3)(r3)(p4), inner sep=0.18cm,
                  label=right:Sync check] {};
        \end{scope}
    \end{tikzpicture}%
    }
    \caption{Hardware and software initialisation sequence that must complete on both laptops before any ROS 2 nodes are launched.}
    \label{fig:ch6_startup_sequence}
\end{figure}

Once the startup sequence completes successfully on both machines, the full collaborative cell is
operational: the station laptop drives the physical hardware and streams telemetry, the control
laptop plans motions and maintains the digital twin, and both channels of the communication
architecture carry their respective payloads across the ZeroTier overlay.